\chapter{2023年北京卷}

\section{单选题}

\begin{problem}
已知集合 \(M=\{x\mid x+2\ge 0\}\)，\(N=\{x\mid x-1<0\}\)，则 \(M\cap N=\)（\quad）

\choices
    {\(\{x\mid -2\le x<1\}\)}
    {\(\{x\mid -2<x\le 1\}\)}
    {\(\{x\mid x\ge -2\}\)}
    {\(\{x\mid x<1\}\)}
\end{problem}

\begin{problem}
在复平面内，复数 \(z\) 对应的点的坐标是 \((-1,\sqrt{3})\)，则 \(z\) 的共轭复数 \(\overline{z}=\)（\quad）

\choices
    {\(1+\sqrt{3}\i\)}
    {\(1-\sqrt{3}\i\)}
    {\(-1+\sqrt{3}\i\)}
    {\(-1-\sqrt{3}\i\)}
\end{problem}

\begin{problem}
已知向量 \(\bs{a}\)，\(\bs{b}\) 满足 \(\bs{a}+\bs{b}=(2,3)\)，\(\bs{a}-\bs{b}=(-2,1)\)，则 \(|\bs{a}|^2-|\bs{b}|^2=\)（\quad）

\choices
    {\(-2\)}
    {\(-1\)}
    {0}
    {1}
\end{problem}

\begin{problem}
下列函数中，在区间 \((0,+\infty)\) 上单调递增的是（\quad）

\choices
    {\(f(x)=-\ln x\)}
    {\(f(x)=\dfrac{1}{2^x}\)}
    {\(f(x)=-\dfrac{1}{x}\)}
    {\(f(x)=3^{|x-1|}\)}
\end{problem}

\begin{problem}
\(\left(2x-\dfrac{1}{x}\right)^5\) 的展开式中 \(x\) 的系数为（\quad）

\choices
    {\(-80\)}
    {\(-40\)}
    {40}
    {80}
\end{problem}

\begin{problem}
已知抛物线 \(C:y^2=8x\) 的焦点为 \(F\)，点 \(M\) 在 \(C\) 上. 若 \(M\) 到直线 \(x=-3\) 的距离为 \(5\)，则 \(|MF|=\)（\quad）

\choices
    {7}
    {6}
    {5}
    {4}
\end{problem}

\begin{problem}
在 \(\triangle ABC\) 中，\((a+c)(\sin A-\sin C)=b(\sin A-\sin B)\)，则 \(\angle C=\)（\quad）

\choices
    {\(\dfrac{\pi}{6}\)}
    {\(\dfrac{\pi}{3}\)}
    {\(\dfrac{2\pi}{3}\)}
    {\(\dfrac{5\pi}{6}\)}
\end{problem}

\begin{problem}
若 \(xy\ne 0\)，则“\(x+y=0\)”是“\(\dfrac{x}{y}+\dfrac{y}{x}=-2\)”的（\quad）

\choices
    {充分不必要条件}
    {必要不充分条件}
    {充要条件}
    {既不充分也不必要条件}
\end{problem}

\begin{problem}
坡屋顶是我国传统建筑造型之一，蕴含着丰富的数学元素. 安装灯带可以勾勒出建筑轮廓，展现造型之美. 如图，某坡屋顶可视为一个五面体，其中两个面是全等的等腰梯形，两个面是全等的等腰三角形. 若 \(AB=25\text{m}\)，\(BC=AD=10\text{m}\)，且等腰梯形所在的平面，等腰三角形所在的平面与平面 \(ABCD\) 的夹角的正切值均为 \(\dfrac{\sqrt{14}}{5}\)，则该五面体的所有棱长之和为（\quad）

\bitmapfigure[max width=6.4cm]{img/2023/beijing/q9_fig1.png}

\choices
    {\(102\text{m}\)}
    {\(112\text{m}\)}
    {\(117\text{m}\)}
    {\(125\text{m}\)}
\end{problem}

\begin{problem}
已知数列 \(\{a_n\}\) 满足 \(a_{n+1}=\frac{1}{4}(a_n-6)^3+6\)（\(n=1\)，\(2\)，\(3\)，\(\dots\)），则（\quad）

\choices
    {当 \(a_1=3\) 时，\(\{a_n\}\) 为递减数列，且存在常数 \(M\le 0\)，使得 \(a_n>M\) 恒成立}
    {当 \(a_1=5\) 时，\(\{a_n\}\) 为递增数列，且存在常数 \(M\le 6\)，使得 \(a_n<M\) 恒成立}
    {当 \(a_1=7\) 时，\(\{a_n\}\) 为递减数列，且存在常数 \(M>6\)，使得 \(a_n>M\) 恒成立}
    {当 \(a_1=9\) 时，\(\{a_n\}\) 为递增数列，且存在常数 \(M>0\)，使得 \(a_n<M\) 恒成立}
\end{problem}

\section{填空题}

\begin{problem}
已知函数 \(f(x)=4^x+\log_2 x\)，则 \(f\left(\dfrac{1}{2}\right)=\fillinblank{}\).
\end{problem}

\begin{problem}
已知双曲线 \(C\) 的焦点为 \((-2,0)\) 和 \((2,0)\)，离心率为 \(\sqrt{2}\)，则 \(C\) 的方程为\fillinblank{}.
\end{problem}

\begin{problem}
已知命题 \(p\)：若 \(\alpha\)，\(\beta\) 为第一象限角，且 \(\alpha>\beta\)，则 \(\tan\alpha>\tan\beta\). 能说明 \(p\) 为假命题的一组 \(\alpha\)，\(\beta\) 的值为 \(\alpha=\fillinblank{}\)，\(\beta=\fillinblank{}\).
\end{problem}

\begin{problem}
我国度量衡的发展有着悠久的历史，战国时期就已经出现了类似于砝码的，用来测量物体质量的“环权”. 已知 9 枚环权的质量（单位：铢）从小到大构成项数为 9 的数列 \(\{a_n\}\)，该数列的前 3 项成等差数列，后 7 项成等比数列，且 \(a_1=1\)，\(a_5=12\)，\(a_9=192\)，则 \(a_7=\fillinblank{}\)；数列 \(\{a_n\}\) 所有项的和为 \(\fillinblank{}\).
\end{problem}

\begin{problem}
设 \(a>0\)，函数 \(f(x)=x+2\)（\(x<-a\)），\(f(x)=\sqrt{a^2-x^2}\)（\(-a\le x\le a\)），\(f(x)=-\sqrt{x}-1\)（\(x>a\)）. 给出下列四个结论：

\begin{circlelist}
\item \(f(x)\) 在区间 \((a-1,+\infty)\) 上单调递减；

\item 当 \(a\ge 1\) 时，\(f(x)\) 存在最大值；

\item 设 \(M(x_1,f(x_1))(x_1\le a)\)，\(N(x_2,f(x_2))(x_2>a)\)，则 \(|MN|>1\)；

\item 设 \(P(x_3,f(x_3))(x_3<-a)\)，\(Q(x_4,f(x_4))(x_4\ge -a)\).
\end{circlelist}
若 \(|PQ|\) 存在最小值，则 \(a\) 的取值范围是 \(\left(0,\dfrac{1}{2}\right]\).

其中所有正确结论的序号是\fillinblank{}.
\end{problem}

\section{解答题}

\begin{problem}
如图，在三棱锥 \(P-ABC\) 中，\(PA\perp \text{平面 }ABC\)，\(PA=AB=BC=1\)，\(PC=\sqrt{3}\).

\begin{enumerate}
\item 求证：\(BC\perp \text{平面 }PAB\)；
\item 求二面角 \(A-PC-B\) 的大小.
\end{enumerate}

\bitmapfigure[max width=4.9cm]{img/2023/beijing/q16_fig1.png}
\end{problem}

\begin{problem}
设函数 \(f(x)=\sin \omega x\cos \varphi+\cos \omega x\sin \varphi\)（\(\omega>0\)，\(|\varphi|<\frac{\pi}{2}\)）.

\begin{enumerate}
\item 若 \(f(0)=-\dfrac{\sqrt{3}}{2}\)，求 \(\varphi\) 的值；
\item 已知 \(f(x)\) 在区间 \(\left[-\dfrac{\pi}{3},\dfrac{2\pi}{3}\right]\) 上单调递增，\(f\left(\dfrac{2\pi}{3}\right)=1\)，再从条件 \(\circled{1}\)，条件 \(\circled{2}\)，条件 \(\circled{3}\) 这三个条件中选择一个作为已知，使函数 \(f(x)\) 存在，求 \(\omega\)，\(\varphi\) 的值.
\end{enumerate}

条件 \(\circled{1}\)：\(f\left(\dfrac{\pi}{3}\right)=\sqrt{2}\)；

条件 \(\circled{2}\)：\(f\left(-\dfrac{\pi}{3}\right)=-1\)；

条件 \(\circled{3}\)：\(f(x)\) 在区间 \(\left[-\dfrac{\pi}{2},-\dfrac{\pi}{3}\right]\) 上单调递减.
\end{problem}

\begin{problem}
为研究某种农产品价格变化的规律，收集得到了该农产品连续 40 天的价格变化数据，如下表所示. 在描述价格变化时，用“+”表示“上涨”，即当天价格比前一天价格高；用“-”表示“下跌”，即当天价格比前一天价格低；用“0”表示“不变”，即当天价格与前一天价格相同. 用频率估计概率.

\begin{center}
\resizebox{\linewidth}{!}{
\begin{tabular}{|c|*{20}{c|}}
\hline
第 1 天到第 20 天 & - & + & + & 0 & - & - & - & + & + & 0 & + & 0 & - & - & + & - & + & 0 & 0 & + \\
\hline
第 21 天到第 40 天 & 0 & + & + & 0 & - & - & - & + & + & 0 & + & 0 & + & - & - & - & + & 0 & - & + \\
\hline
\end{tabular}
}
\end{center}

\begin{enumerate}
\item 试估计该农产品价格“上涨”的概率；
\item 假设该农产品每天的价格变化是相互独立的. 在未来的日子里任取 4 天，试估计该农产品价格在这 4 天中 2 天“上涨”，1 天“下跌”，1 天“不变”的概率；
\item 假设该农产品每天的价格变化只受前一天价格变化的影响. 判断第 41 天该农产品价格“上涨”“下跌”和“不变”的概率估计值哪个最大. （结论不要求证明）
\end{enumerate}
\end{problem}

\begin{problem}
已知椭圆 \(E:\frac{x^2}{a^2}+\frac{y^2}{b^2}=1\)（\(a>b>0\)） 的离心率为 \(\dfrac{\sqrt{5}}{3}\)，\(A\)，\(C\) 分别是 \(E\) 的上，下顶点，\(B\)，\(D\) 分别是 \(E\) 的左，右顶点，\(|AC|=4\).

\begin{enumerate}
\item 求 \(E\) 的方程；
\item 设 \(P\) 为第一象限内 \(E\) 上的动点，直线 \(PD\) 与直线 \(BC\) 交于点 \(M\)，直线 \(PA\) 与直线 \(y=-2\) 交于点 \(N\). 求证：\(MN\parallel CD\).
\end{enumerate}
\end{problem}

\begin{problem}
设函数 \(f(x)=x-x^3\e^{ax+b}\)，曲线 \(y=f(x)\) 在点 \((1,f(1))\) 处的切线方程为 \(y=-x+1\).

\begin{enumerate}
\item 求 \(a\)，\(b\) 的值；
\item 设函数 \(g(x)=f'(x)\)，求 \(g(x)\) 的单调区间；
\item 求 \(f(x)\) 的极值点个数.
\end{enumerate}
\end{problem}

\begin{problem}
已知数列 \(\{a_n\}\)，\(\{b_n\}\) 的项数均为 \(m\)（\(m>2\)），且 \(a_n\)，\(b_n\in\{1,2,\dots,m\}\)，\(\{a_n\}\)，\(\{b_n\}\) 的前 \(n\) 项和分别为 \(A_n\)，\(B_n\)，并规定 \(A_0=B_0=0\). 对于 \(k\in\{0,1,2,\dots,m\}\)，定义 \(r_k=\max \{ i\mid B_i\le A_k,\ i\in\{0,1,2,\dots,m\}\}\)，其中，\(\max M\) 表示数集 \(M\) 中最大的数.

\begin{enumerate}
\item 若 \(a_1=2\)，\(a_2=1\)，\(a_3=3\)，\(b_1=1\)，\(b_2=3\)，\(b_3=3\)，求 \(r_0\)，\(r_1\)，\(r_2\)，\(r_3\) 的值；
\item 若 \(a_1\ge b_1\)，且 \(2r_j\le r_{j+1}+r_{j-1}\)（\(j=1\)，\(2\)，\(\dots\)，\(m-1\)），求 \(r_n\)；
\item 证明：存在 \(p\)，\(q\)，\(s\)，\(t\in\{0,1,2,\dots,m\}\)，满足 \(p>q\)，\(s>t\)，使得 \(A_p+B_t=A_q+B_s\).
\end{enumerate}
\end{problem}

